% Complete proof of C_F0068 Part A, Sections 2--3; no Part B example.
\section{Rationality under maximal \texorpdfstring{$\Delta$}{Delta}}\label{sec:maximal-delta}

We prove Theorem~\ref{thm:delta} in this section.
By Definition~\ref{def:psi}, $\Psi_i$ is right multiplication by $X_i\otimes1-1\otimes X_i^{\mathrm{op}}$.
In particular,
\begin{equation}\label{eq:psi-left-linearity}
 \Psi_i(rk)=r\Psi_i(k),\qquad
 r,k\in L^0(\M)\otimes_{\C}L^0(\M)^{\mathrm{op}}.
\end{equation}

\begin{lemma}\label{lem:tensor-clearing}
Given $k_1,\ldots,k_d\in L^0(\M)\otimes_{\C}L^0(\M)^{\mathrm{op}}$, there are regular $s,t\in\M$ such that
\[
 (s\otimes t^{\mathrm{op}})k_i
 \in\M\otimes_{\C}\M^{\mathrm{op}}\qquad(1\leq i\leq d).
\]
\end{lemma}
\begin{proof}
For simplicity of notation, we give the calculation for simple tensors
\[
 k_i=u_i\otimes v_i^{\mathrm{op}},\qquad u_i,v_i\in L^0(\M).
\]
The same argument applies to general finite sums.
Applying Lemma~\ref{lem:clearing} to $u_1,\ldots,u_d$ gives a regular $s\in\M$ such that
\[
 su_i\in\M\qquad(1\leq i\leq d).
\]
Applying its right-denominator assertion to $v_1,\ldots,v_d$ gives a regular $t\in\M$ such that
\[
 v_it\in\M\qquad(1\leq i\leq d).
\]
Then
\[
 (s\otimes t^{\mathrm{op}})k_i
 =su_i\otimes(v_it)^{\mathrm{op}}
 \in\M\otimes_{\C}\M^{\mathrm{op}}.
\]
\end{proof}

\begin{lemma}\label{lem:tensor-injectivity}
Suppose $\Delta(X)=d$.
The map
\[
 (L^0(\M)\otimes_{\C}L^0(\M)^{\mathrm{op}})^d\longrightarrow
 L^0(\M)\otimes_{\C}L^0(\M)^{\mathrm{op}},\qquad
 (k_i)\longmapsto\sum_i\Psi_i(k_i)
\]
is injective.
\end{lemma}
\begin{proof}
Suppose $\sum_i\Psi_i(k_i)=0$.
Choose regular $s,t\in\M$ as in Lemma~\ref{lem:tensor-clearing} and set $T_i:=\Theta((s\otimes t^{\mathrm{op}})k_i)\in\F(H)$.
By \eqref{eq:tensor-commutator} and \eqref{eq:psi-left-linearity},
\[
 \sum_i[\rho(X_i),T_i]
 =\Theta\left((s\otimes t^{\mathrm{op}})\sum_i\Psi_i(k_i)\right)=0.
\]
Remark~\ref{rem:cycle-characterization} gives $T_i=0$ for every $i$.
Injectivity of $\Theta$ implies $(s\otimes t^{\mathrm{op}})k_i=0$.
Left multiplication by $s^{-1}\otimes(t^{-1})^{\mathrm{op}}$ then gives $k_i=0$.
\end{proof}

\begin{lemma}\label{lem:tensor-certificate}
Let $Q=(s,a,b,t)$ be a fraction pair for $u\in L^0(\M)$, and let $\ell_1,\ldots,\ell_d\in H\otimes_{\C}H$.
Set $T_i:=-\Theta(\ell_i)\in\F(H)$.
Then
\[
 \sum_i\Psi_i(\ell_i)
 =a\otimes t^{\mathrm{op}}-s\otimes b^{\mathrm{op}}
\]
if and only if
\[
 \sum_i[\rho(X_i),T_i]=\Gamma[Q,P_\Omega].
\]
\end{lemma}
\begin{proof}
By \eqref{eq:tensor-commutator} and \eqref{eq:theta-bounded},
\[
 \sum_i[\rho(X_i),T_i]
 =-\Theta\left(\sum_i\Psi_i(\ell_i)\right),
 \qquad
 \Gamma[Q,P_\Omega]
 =-\Theta(a\otimes t^{\mathrm{op}}-s\otimes b^{\mathrm{op}}).
\]
The equivalence follows from the injectivity of $\Theta$.
\end{proof}

By Lemma~\ref{lem:free-field}, $\Div_{L^0(\M)}\A_X$ is a division subring of $L^0(\M)$ whenever $\Delta(X)=d$.
We prove Theorem~\ref{thm:delta} through $(1)\implies(2)\implies(3)\implies(4)\implies(5)\implies(1)$.

\begin{proof}[Proof of Theorem~\ref{thm:delta}]
In this proof, set $U:=L^0(\M)$ and $F:=\Div_U\A_X$.

$(1)\implies(2)$.
By \cite[Theorem~1.1, (i) and (v)]{MSY23}, the assumption $\Delta(X)=d$ implies that $X_1,\ldots,X_d$ generate the free field: evaluation identifies $\freefield{x_1,\ldots,x_d}$ with $F$.
The free difference quotients from Subsection~\ref{subsec:free-field}, transported through this evaluation isomorphism, define linear maps
\[
 \partial_i:F\longrightarrow F\otimes_{\C}F^{\mathrm{op}}.
\]
They satisfy the rules in equation~\eqref{eq:free-field-difference-rules} with $x_j$ replaced by $X_j$.
Put $\delta v:=v\otimes1-1\otimes v^{\mathrm{op}}$.
Direct expansion gives
\[
 \delta(v_1v_2)=(1\otimes v_2^{\mathrm{op}})\delta v_1+(v_1\otimes1)\delta v_2,
 \qquad
 \delta(v^{-1})=-(v^{-1}\otimes(v^{-1})^{\mathrm{op}})\delta v.
\]
For constants, $\delta c=0$, while for generators, $\delta X_j=\Psi_j(1\otimes1)=\sum_i\Psi_i(\partial_iX_j)$.
We now show that this identity extends from constants and generators to all rational elements using the product and inverse rules.

If $\delta v_\ell=\sum_i\Psi_i(\partial_i v_\ell)$ for $\ell=1,2$, the product rules and equation~\eqref{eq:psi-left-linearity} give
\begin{align*}
 \delta(v_1v_2)
 &=(1\otimes v_2^{\mathrm{op}})\sum_i\Psi_i(\partial_i v_1)
   +(v_1\otimes1)\sum_i\Psi_i(\partial_i v_2)\\
 &=\sum_i\Psi_i\bigl((1\otimes v_2^{\mathrm{op}})\partial_i v_1
                  +(v_1\otimes1)\partial_i v_2\bigr)\\
 &=\sum_i\Psi_i\bigl(\partial_i(v_1v_2)\bigr).
\end{align*}
Similarly, if $0\ne v\in F$ satisfies $\delta v=\sum_i\Psi_i(\partial_i v)$, the inverse rules give
\begin{align*}
 \delta(v^{-1})
 &=-(v^{-1}\otimes(v^{-1})^{\mathrm{op}})\sum_i\Psi_i(\partial_i v)\\
 &=\sum_i\Psi_i\bigl(-(v^{-1}\otimes(v^{-1})^{\mathrm{op}})\partial_i v\bigr)\\
 &=\sum_i\Psi_i\bigl(\partial_i(v^{-1})\bigr).
\end{align*}
Together with linearity, this shows that the elements satisfying the identity form a division subalgebra of $F$ containing $\C$ and $X_1,\ldots,X_d$.
Since these generate $F$, the identity holds for every rational element $u\in F$.
This proves equation~\eqref{eq:tensor-criterion} with $k_i:=\partial_i u$; thus assertion~\textup{(\ref{item:delta-tensors})} holds.
Lemma~\ref{lem:tensor-injectivity} proves uniqueness of the tensors in assertion~\textup{(\ref{item:delta-tensors})}.

$(2)\implies(3)$.
Given tensors $k_i$ as in assertion~\textup{(2)}, for any fraction pair $Q=(s,a,b,t)$ set $\ell_i:=(s\otimes t^{\mathrm{op}})k_i$.
By equation~\eqref{eq:psi-left-linearity},
\[
 \sum_i\Psi_i(\ell_i)
 =(s\otimes t^{\mathrm{op}})
   (u\otimes1-1\otimes u^{\mathrm{op}})
 =a\otimes t^{\mathrm{op}}-s\otimes b^{\mathrm{op}},
\]
so assertion~\textup{(\ref{item:delta-all-fraction-tensors})} holds.

$(3)\implies(4)$.
Choose a fraction pair $Q=(s,a,b,t)$ for $u$ and corresponding tensors $\ell_i$ in assertion~\textup{(3)}.
By Lemma~\ref{lem:tensor-clearing}, there are regular $s',t'\in\M$ such that
\[
 (s'\otimes(t')^{\mathrm{op}})\ell_i
 \in\M\otimes_{\C}\M^{\mathrm{op}}.
\]
By equation~\eqref{eq:psi-left-linearity} and equation~\eqref{eq:all-fraction-tensor-criterion},
\[
 \sum_i\Psi_i\bigl((s'\otimes(t')^{\mathrm{op}})\ell_i\bigr)
 =s'a\otimes(tt')^{\mathrm{op}}-s's\otimes(bt')^{\mathrm{op}}.
\]
Since $(s's,s'a,bt',tt')$ is a fraction pair for $u$, assertion~\textup{(\ref{item:delta-bounded-fraction-tensors})} holds.

For fixed $Q$, uniqueness in assertions \textup{(\ref{item:delta-all-fraction-tensors})} and \textup{(\ref{item:delta-bounded-fraction-tensors})} follows from Lemma~\ref{lem:tensor-injectivity}, applied to the difference of two solutions of equation~\eqref{eq:all-fraction-tensor-criterion}.

$(4)\implies(5)$.
Let $Q$ and $\ell_i$ be as in assertion~\textup{(\ref{item:delta-bounded-fraction-tensors})}.
Lemma~\ref{lem:tensor-certificate} gives the finite-rank operators $T_i:=-\Theta(\ell_i)$ satisfying equation~\eqref{eq:certificate}.
Thus assertion~\textup{(\ref{item:delta-finite-rank})} holds.

$(5)\implies(1)$.
Suppose assertion~\textup{(\ref{item:delta-finite-rank})} holds.
Let $Q=(s,a,b,t)$ be the fraction pair and $T_1,\ldots,T_d\in\F(H)$ the operators satisfying equation~\eqref{eq:certificate}.
Via $\Theta$, define a linear map $\pi:\F(H)\longrightarrow U\otimes_FU$ by
\[
 \pi(\Theta(v_1\otimes v_2^{\mathrm{op}})):=v_1\otimes_Fv_2,
 \qquad v_1,v_2\in H.
\]
Note that $\Theta$ is a linear bijection, so $\pi$ is well defined.
For a rank-one operator $\Theta(v_1\otimes v_2^{\mathrm{op}})$, equation~\eqref{eq:tensor-commutator} gives
\begin{align*}
 \pi\bigl([\rho(X_i),\Theta(v_1\otimes v_2^{\mathrm{op}})]\bigr)
 &=(v_1X_i)\otimes_Fv_2-v_1\otimes_F(X_iv_2)\\
 &=0.
\end{align*}
The last equality uses $X_i\in F$ and the balancing relation
\[
 (uf)\otimes_Fw=u\otimes_F(fw),\qquad u,w\in U,\quad f\in F.
\]
Every finite-rank operator is a finite sum of rank-one operators, so linearity gives $\pi([\rho(X_i),T])=0$ for every $T\in\F(H)$.

Applying $\pi$ to equation~\eqref{eq:certificate} and using $\Theta(v_1\otimes v_2^{\mathrm{op}})=v_1P_\Omega v_2$ for $v_1,v_2\in\M$ (see equation~\eqref{eq:theta-bounded}), we obtain
\begin{align*}
 0&=\pi\left(\sum_i[\rho(X_i),T_i]\right)\\
  &=\pi(\Gamma[Q,P_\Omega])\\
  &=\pi(sP_\Omega b-aP_\Omega t)\\
  &=\pi\bigl(\Theta(s\otimes b^{\mathrm{op}})-\Theta(a\otimes t^{\mathrm{op}})\bigr)\\
  &=s\otimes_Fb-a\otimes_Ft.
\end{align*}
Hence $a\otimes_Ft=s\otimes_Fb$.

To clear the denominators, consider the map
\[
 \Phi_{s,t}:U\otimes_FU\longrightarrow U\otimes_FU,\qquad
 \Phi_{s,t}(u\otimes_Fw):=s^{-1}u\otimes_Fwt^{-1}.
\]
It is well defined because it preserves the balancing relation.
Applying $\Phi_{s,t}$ to $a\otimes_Ft=s\otimes_Fb$ yields
\[
 s^{-1}a\otimes_Ftt^{-1}=s^{-1}s\otimes_Fbt^{-1}.
\]
Since $u=s^{-1}a=bt^{-1}$, this is
\[
 u\otimes_F1=1\otimes_Fu.
\]
Lemma~\ref{lem:equalizer} therefore gives $u\in F=\Div_U\A_X$, so assertion~\textup{(1)} holds.

For fixed $Q$, let $(T_1,\ldots,T_d)$ and $(S_1,\ldots,S_d)$ be two finite-rank solutions of equation~\eqref{eq:certificate}.
Their right-hand sides are both $\Gamma[Q,P_\Omega]$, so subtraction gives
\begin{align*}
 \sum_i[\rho(X_i),T_i-S_i]
 &=\sum_i[\rho(X_i),T_i]-\sum_i[\rho(X_i),S_i]=0.
\end{align*}
Each $T_i-S_i$ has finite rank.
Since $\Delta(X)=d$, Remark~\ref{rem:cycle-characterization} implies $T_i-S_i=0$ for every $i$.
Hence the two tuples coincide, proving uniqueness of the finite-rank solution tuple for the fixed fraction pair $Q$.

\end{proof}

\begin{remark}\label{rem:boundary-interpretation}
Equip $L^0(\M)\otimes_{\C}L^0(\M)^{\mathrm{op}}$ with the $L^0(\M)$-bimodule structure defined, for $u,v\in L^0(\M)$ and $k\in L^0(\M)\otimes_{\C}L^0(\M)^{\mathrm{op}}$, by
\[
 u\cdot k\cdot v:=k(v\otimes u^{\mathrm{op}}).
\]
The degree-one Hochschild differential with these coefficients \cite[Section~9.1.1, p.~300]{Wei94} is the map
\[
 d_1:\bigl(L^0(\M)\otimes_{\C}L^0(\M)^{\mathrm{op}}\bigr)
 \otimes_{\C}L^0(\M)
 \longrightarrow L^0(\M)\otimes_{\C}L^0(\M)^{\mathrm{op}}
\]
given by
\[
 d_1(k\otimes u):=k\cdot u-u\cdot k
 =k(u\otimes1-1\otimes u^{\mathrm{op}}).
\]
Put $V_X:=\operatorname{span}_{\C}\{X_1,\ldots,X_d\}$.
Restricting the last tensor factor to $V_X$, we obtain
\[
 \partial_X:
 \bigl(L^0(\M)\otimes_{\C}L^0(\M)^{\mathrm{op}}\bigr)
 \otimes_{\C}V_X
 \longrightarrow L^0(\M)\otimes_{\C}L^0(\M)^{\mathrm{op}},
 \qquad \partial_X:=d_1|_{(L^0(\M)\otimes_{\C}L^0(\M)^{\mathrm{op}})\otimes_{\C}V_X}.
\]
By Definition~\ref{def:psi},
\[
 \partial_X\!\left(\sum_i k_i\otimes X_i\right)
 =\sum_i\Psi_i(k_i).
\]
This restriction to the linear span of the generators follows \cite[Section~3.1.1, equation~(3.1)]{CS05}, with $L^0(\M)\otimes_{\C}L^0(\M)^{\mathrm{op}}$ as the coefficient bimodule.
Thus condition~\textup{(\ref{item:delta-tensors})} in Theorem~\ref{thm:delta} requires $u\otimes1-1\otimes u^{\mathrm{op}}$ to lie in the image of $\partial_X$.

Similarly, equip $\F(H)$ with the $\A_X$-bimodule structure
\[
 x\cdot T\cdot y:=\rho(y)T\rho(x),
 \qquad x,y\in\A_X,\quad T\in\F(H).
\]
The degree-one Hochschild differential is then
\[
 d_1:\F(H)\otimes_{\C}\A_X\longrightarrow\F(H),
 \qquad d_1(T\otimes x):=T\cdot x-x\cdot T=[\rho(x),T].
\]
Since $\A_X$ is generated by $X_1,\ldots,X_d$, the identity
\[
 [\rho(xy),T]
 =[\rho(y),\rho(x)T]+[\rho(x),T\rho(y)]
\]
shows, by induction on monomials, that restricting the last tensor factor to $V_X$ gives the full image of $d_1$:
\[
 \operatorname{im}d_1
 =d_1\bigl(\F(H)\otimes_{\C}V_X\bigr)
 =\left\{\sum_{i=1}^d[\rho(X_i),T_i]:T_i\in\F(H)\right\}.
\]
Condition \eqref{eq:certificate} is therefore
\[
 d_1\!\left(\sum_iT_i\otimes X_i\right)=\Gamma[Q,P_\Omega].
\]
Thus $\Gamma[Q,P_\Omega]$ is a boundary, or equivalently its class vanishes in $H_0(\A_X,\F(H)):=\F(H)/\operatorname{im}d_1$.
\end{remark}
